\begin{textsample}{POS Dim 2 – vem\_ed\_19 – Score 14.00 – vem\_ed\_19/t091.txt}  \label{ex:f2_pos_009}
continuação Alguns sinais de que um projeto não vai bem: professores que não participam, não \textbf{detalham} o projeto; docentes que culpam o parceiro ou não buscam soluções / bons resultados para o projeto - considerado como atividade \textbf{pontual} sem \textbf{continuidade} e sem visão das relações \textbf{interinstitucionais}; professores que não \textbf{escutam} seus alunos, não seguem o \textbf{planejado} ou mudam os planos no meio do projeto sem consultar parceiros e não buscam um “\textbf{justo} \textbf{meio-termo}”; estudantes que não interagem, confusos ou em constante conflito, sem alguém para coordená-los, ou ainda insensíveis à cultura do outro país, vista sem respeito e com \textbf{estereótipos}.
Do ponto de vista institucional, portanto, é imprescindível contar com um coordenador que ajude a resolver os problemas locais.
Em resumo: para expandir os Intercâmbios Virtuais com \textbf{qualidade}, é \textbf{preciso} ter uma equipe que \textbf{gerencie} o \textbf{design}, a \textbf{implantação} e o desenvolvimento desses projetos.

% matched lemmas: continuidade, design, detalhar, escutar, estereótipo, gerenciar, implantação, interinstitucional, justo, meio-termo, planejado, pontual, preciso, qualidade
\end{textsample}
